\documentclass[11pt]{article}
\usepackage[margin=1in]{geometry}
\usepackage{times}
\usepackage{hyperref}
\hypersetup{colorlinks=true, linkcolor=black, urlcolor=blue, citecolor=black}
\title{Rotating a Magnet Faster Does Not Rotate the Rule Book: Paramahamsa Tewari's Space Power Generator}
\author{Flyxion \\ \small Independent Researcher}
\date{}
\begin{document}
\maketitle

\begin{abstract}
Paramahamsa Tewari, a former Indian nuclear engineer, has since the 1980s promoted a Space Power Generator that he claimed could deliver electrical output exceeding the mechanical power required to drive it by exploiting a rotor and coil geometry said to draw energy directly from the surrounding vacuum or from a medium he calls space. The claim is tied to a broader theoretical framework in which matter and charge arise from a substratum underlying standard particle physics. This paper argues that the reported efficiency gains trace to a common measurement error in generator claims: the drive motor's electrical input is treated as if it were the true mechanical shaft input, without the measured torque and angular speed that determine the real mechanical work delivered. It also argues that no published, independent, shaft-power test in the literature I have been able to verify establishes the claimed excess under the standards required for such a claim. Tewari's broader theory, while internally elaborate, has not been shown to reproduce the predictive successes of standard electromagnetic theory in the domains where that theory is already highly confirmed. The paper therefore treats the device as an example of a claim that may be physically interesting but has not established itself as a demonstration of new energy physics.
\end{abstract}

\section{Introduction}

Paramahamsa Tewari, who worked as a senior engineer in India's nuclear sector, has since the 1980s promoted a generator design that he has described as a space-based power source, a machine in which a rotating magnet and stationary coil geometry are said to yield more electrical output than the real mechanical input needed to sustain the rotation. Tewari's case is different from some other over-unity proposals in that he has sometimes framed the device within an explicit theoretical program. In his account, the ordinary electromagnetic field is not the ultimate structure of reality but the emergent expression of a deeper medium he refers to as space or absolute space. Mass and charge are made out of this underlying structure rather than simply being properties of ordinary particles in a vacuum. The generator is then treated not merely as an unusual machine but as evidence for a more general theory of physical reality.

The machine itself is a generator claim, not a direct theory of the cosmos. It is therefore best assessed by the standards of electrical engineering: how much mechanical torque is delivered to the shaft, what the rotational speed is, what the electrical output is, and how much of that output is actually available after internal losses. The power balance cannot be judged by the motor's electrical input alone. A motor may be supplied with power from a source, but the shaft input to the generator is not determined by the source reading unless the motor efficiency, coupling efficiency, transmission losses, and actual operating point are known. This distinction is the key to why so many generator claims look impressive in the first presentation and collapse under a shaft-power audit.

The relevant historical question is therefore not whether Tewari's background or his theoretical framework is sophisticated. The relevant question is whether the mechanical-to-electrical conversion has been measured at the level of rigor that the claim requires. In the case of the Space Power Generator, the published evidence I have been able to verify does not resolve that question in the way an over-unity claim requires. The claim remains a contested device story, not a demonstrated exception to the energy balance.

\section{The Standard Error in This Class of Claim}

The recurring error in this class of claim is not exotic. It is that the input power to the drive motor is treated as if it were the actual mechanical power transmitted to the generator shaft. A motor's electrical input is not the same as the shaft power it delivers, because some energy is lost in the conversion from electrical to mechanical form. The mismatch becomes especially large if the motor is being run off a source whose source rating is not the same as the power transmitted at the shaft, or if the coupling or gearbox has efficiency losses, or if the generator's own drag and brush losses change with load. If the generator is loaded differently from a conventional machine, the apparent efficiency can easily become inflated when the wrong input quantity is used.

That issue is not a footnote. It is exactly what turns an ordinary machine into a supposed over-unity outlier. A generator can appear to exceed one hundred percent efficiency when the drive-motor electrical input is used as the input measure while the actual shaft work is not directly measured. This is why a direct torque-sensor test is necessary. The standard quantity for ordinary generator evaluation is not the drive-motor power reading. It is the mechanical torque at the shaft multiplied by the angular speed. Only then can the generated electrical power be compared to the real mechanical input under the same operating conditions.

This is the crucial point in Tewari's case. If the device is said to produce more power than the mechanical input, then the input must be measured at the shaft, not inferred from source power or motor input. A design that reduces magnetic drag or back reaction can be genuinely interesting, but it does not change the fact that mechanical work is still the source of the electrical output. A lower-drag machine is not a free-energy machine. It is simply a machine with lower parasitic losses. The claim requires more than a lower measured drag. It requires an independent shaft-power measurement showing that the output exceeds the actual shaft work delivered under controlled conditions.

\section{The Main Theoretical Question Is Not Whether the Machine Is Unusual}

Tewari's generator is unusual in the way many over-unity designs are unusual: it is built around a claim that certain rotor and stator arrangements change the mechanical consequences of generating current. In a conventional generator, as current is drawn, the field exerts a magnetic drag torque on the shaft. This is the ordinary back-reaction of electromagnetic induction. A designer can reduce drag by changing the geometry, by improving the magnet arrangement, or by altering the current paths. These are real engineering considerations. The existence of lower drag does not in itself show that the machine is extracting energy from a hidden source. It shows only that the machine has been optimized in one dimension of the energy conversion problem.

This is the distinction that matters. A claim that a specific design produces unusual low back torque, or reduced magnetic drag for a given output, is not a claim that the machine has broken the law of energy conservation. It is merely a claim that the field or geometry has been shaped in a different way. A real over-unity claim would require direct evidence that the machine produces more electrical energy than the shaft work it receives, measured at the shaft, under controlled conditions. That is a higher standard and one the publicly available records of this generator have not clearly met.

The same issue appears in other devices in this volume, especially the N-machine claim and other generator designs that are evaluated by source power rather than shaft power. This pattern is not a mere coincidence. It is a common way in which a machine can look anomalous before the mechanical power is measured in the correct place. The crucial question is therefore not whether a rotor geometry is unusual enough to sound novel. It is whether the output exceeds the real mechanical input after the load, losses, and torque relationship are all measured properly.

\section{The Theoretical Framework Does Not Solve the Measurement Problem}

Tewari's broader theory is often cited by supporters as a reason to take the generator seriously. He has proposed that the underlying vacuum or space medium is not a passive background but a structured field in which matter and charge emerge from more fundamental configurations. This is not an empty point. It matters because many unusual energy claims are tied to unusual metaphysical frameworks. The problem is that a deep theory does not itself validate an engineering claim. A replacement theory must, at minimum, reproduce the empirical successes of the standard theory that it seeks to replace. It must explain ordinary electromagnetic behavior, not merely propose an alternative vocabulary for it.

This is exactly the stringent bar that Tewari's framework has not been shown to clear. Standard electromagnetic theory is not a vague set of intuitions. It is a predictive framework that explains a vast array of phenomena, from generator operation to spectroscopy to material response, with extraordinary precision. A claim that the underlying medium is responsible for the effect is not enough. The theory must show that the actual electromagnetic behavior in the relevant regime follows from the theory. Without that demonstration, the theory has not moved from speculation to support. The same logic applies elsewhere in the series: a novel theoretical picture can be interesting, but it cannot rescue a claim of excess energy if the engineering measurement is still unresolved.

\section{The Entropic Gravity Framing}

The Tewari generator is not primarily a gravitational claim, so the entropic-gravity framing here is brief. In the broader research program behind this series, gravitational curvature is treated as the large-scale expression of an ordered configuration rather than as a field that can be locally filtered or shielded by an ordinary rotating machine in the way an electromagnetic field can be. A generator's rotating magnet may change the electromagnetic conditions of a machine, but it does not automatically show that gravity or spacetime structure has been altered in a way that explains the power claim. The relevant problem remains electromechanical, not gravitational. A machine still has to satisfy the ordinary torque-and-speed energy balance at the shaft before any more exotic theoretical account can carry weight.

\section{Objections}

The most serious defense of the generator is that Tewari's design does not simply mimic a standard generator but changes the way the magnetic field interacts with the rotor and the current path, and therefore it should not be judged by the same crude metric as a conventional machine. This is a legitimate objection in the sense that a genuinely novel generator geometry may change the loss profile. But it is not a reason to suspend the shaft-power test. Novel geometry is not evidence of a new source of energy. It is evidence that the geometry may have a different loss distribution. Unless direct torque measurement shows that the electrical output exceeds the real shaft input, the claim remains in the domain of design curiosity rather than verified excess power.

A second objection is that a professional engineer's background gives the theory and the machine a degree of credibility that a casual inventor's claim would not have. This is not enough. Professional credentials do not substitute for independent, controlled measurement. An engineer can design a machine with a measurement error, and the value of the engineering background does not rescue a power claim when the wrong quantity has been used as the input. The correct question remains the same: what are the actual torque and speed at the shaft, and how much electrical power is produced under those conditions?

\section{Conclusion}

The Space Power Generator is best understood as a cautionary example of a common generator failure mode: the use of a convenient electrical proxy in place of the true shaft input. The reported output may look extraordinary when expressed against motor electrical input or source power, but that does not determine the actual power conversion efficiency of the machine. The real test is the shaft torque and angular speed, measured directly, and the actual output compared to that mechanical work. Tewari's broader theoretical framework, while ambitious, does not by itself establish the generator claim, because it has not been shown to reproduce the standard electromagnetic predictions that the claim depends upon. The result is therefore not a confirmed over-unity machine but an unverified claim whose power accounting remains to be established under the instrumentation standards required for the claim.

\begin{thebibliography}{99}

\bibitem{Tewari2021} P. Tewari, \emph{Space Power Generator}, 2021. 

\bibitem{Tewari2010} P. Tewari, \emph{The Theory of Space Energy}, 2010.

\bibitem{Jackson1999} J. D. Jackson, \emph{Classical Electrodynamics}, 3rd ed., Wiley, 1999.

\end{thebibliography}

\end{document}
